\documentclass[12pt,a4paper]{article}
\usepackage{goyabase}

\usepackage{longtable}
\usepackage{calc}

% Pandoc compatibility
\providecommand{\tightlist}{%
  \setlength{\itemsep}{0pt}\setlength{\parskip}{0pt}}
\newcommand{\passthrough}[1]{#1}

\title{\textbf{Redis en Docker - UT8 NoSQL}}
\author{Departamento de Informática}
\date{\today}

\usepackage[spanish, es-noshorthands, es-tabla]{babel}
\usepackage[autostyle]{csquotes}
\begin{document}

\maketitle

\hypertarget{redis-en-docker---ut8-nosql}{%
\section{Redis en Docker - UT8
NoSQL}\label{redis-en-docker---ut8-nosql}}

Este entorno permite trabajar con \textbf{Redis}, una base de datos NoSQL de tipo \textbf{Clave-Valor}.

\hypertarget{contenidos}{%
\subsection{Contenidos}\label{contenidos}}

\begin{enumerate}
\def\labelenumi{\arabic{enumi}.}
\tightlist
\item
  \textbf{Redis Server}: El motor de base de datos en el puerto
  \texttt{6379}.
\item
  \textbf{RedisInsight}: Interfaz gráfica para visualizar los datos en
  el puerto \texttt{5540}.
\item
  \textbf{Directorio \texttt{src}}: Contiene
  \texttt{datos\_iniciales.redis} con +1000 registros
  para prácticas de carga masiva.
\end{enumerate}

\hypertarget{instrucciones-de-uso}{%
\subsection{Instrucciones de uso}\label{instrucciones-de-uso}}

\subsubsection{Configuración del Entorno (Docker Compose)}\label{configuracion-del-entorno}

Este archivo define los servicios necesarios para la práctica:

\begin{lstlisting}[language=YAML]
services:
  redis:
    image: redis:latest
    container_name: redis-server
    ports:
      - "6379:6379"
    volumes:
      - redis_data:/data
      - ./src:/src
    restart: always

  redisinsight:
    image: redis/redisinsight:latest
    container_name: redisinsight
    ports:
      - "5540:5540"
    depends_on:
      - redis
    restart: always

volumes:
  redis_data:
\end{lstlisting}

\subsubsection{Puesta en marcha}\label{puesta-en-marcha}

Para arrancar el servicio:

\begin{lstlisting}[language=bash]
docker compose up -d
\end{lstlisting}

Para realizar la carga masiva del catálogo inicial (que necesitarás para
la práctica):

\begin{lstlisting}[language=bash]
docker exec -i redis-server redis-cli < src/datos_iniciales.redis
\end{lstlisting}

Para acceder a la consola de Redis (CLI):

\begin{lstlisting}[language=bash]
docker exec -it redis-server redis-cli
\end{lstlisting}

Para acceder a la interfaz gráfica: Abre tu navegador en
\texttt{http://localhost:5540}.

\hypertarget{comandos-buxe1sicos-para-practicar}{%
\subsection{Comandos Básicos para
practicar}\label{comandos-buxe1sicos-para-practicar}}

\begin{lstlisting}
SET user:105:nombre "Víctor"
GET user:105:nombre
INCR metricas:visitas:inicio
LPUSH busquedas:recientes "Paris"
LPUSH busquedas:recientes "Roma"
LRANGE busquedas:recientes 0 -1
\end{lstlisting}

\section{Ejemplos de uso}

Consulta los ejemplos de uso de la empresa Globotour en \texttt{01\_Ejemplos\_Uso.pdf}.

\end{document}
